\begin{textsample}{POS Dim 1 – ac – Score 43.00 – ac/ac\_ef\_46.txt}  \label{ex:f1_pos_180}
4.2.
Anos \textbf{finais} O desafio em relação ao trabalho com a Língua Portuguesa, passados os primeiros anos de escolaridade, é consolidar e garantir a continuidade do que foi aprendido e a superação de dificuldades que eventualmente se tenham acumulado.
Para \textbf{tanto}, é necessário que se investigue que conhecimentos já foram adquiridos e consolidados pelos \textbf{alunos}, a fim de planejar as intervenções pedagógicas adequadas para assegurar que os \textbf{alunos} continuem aprendendo, visando ainda o nível de complexidade que se pretende alcançar.
Assim, nos anos posteriores do Ensino Fundamental, a \textbf{tarefa} \textbf{central} é aprofundar possibilidades de uso da linguagem oral e escrita e dos conhecimentos linguísticos que contribuem para \textbf{tanto}, oferecendo condições para que cada \textbf{aluno} possa desenvolver cada vez mais sua autonomia como usuário da língua.
No que \textbf{diz} respeito à ampliação dos conteúdos, tem -se a exploração das práticas de linguagem \textbf{contemporâneas}, que estão relacionadas não só aos gêneros \textbf{digitais}, com conteúdo cada vez mais \textbf{multissemiótico} ( texto com a junção de diversas linguagens, por exemplo, o uso de recursos verbais e imagéticos ) e multimidiáticos ( veiculados e produzidos em diferentes mídias, áudio, \textbf{vídeo} etc. ), como também às novas formas de produzir, interagir, veicular e \textbf{configurar}.
À escola, cabe contemplar essas novas práticas de produção, \textbf{tanto} no que \textbf{diz} respeito ao uso de modo ético, quanto \textbf{fomentar} \textbf{debates} acerca de temas que circundam essas atividades.
O objetivo é conduzir o \textbf{aluno} para o domínio não apenas dos gêneros escritos abordados pela escola há bastante tempo, mas também proporcionar -lhe o desenvolvimento dos novos \textbf{letramentos}.
Desse modo, é importante destacar que toda a estrutura do currículo de Língua Portuguesa dialoga com as dez \textbf{competências} gerais da Educação Básica, de \textbf{maneira} articulada, a fim de promover a construção dos conhecimentos necessários ao desenvolvimento das habilidades, bem como atitudes e valores, por meio das \textbf{competências} específicas da área de Linguagens.
Assim, a expectativa é de que os \textbf{alunos} sejam capazes de : 6º ano - Utilizar a linguagem oral de forma adequada, em diferentes situações comunicativas, respeitando os diferentes modos de falar. - Utilizar, com propriedade, os conhecimentos sobre padrões da escrita \textbf{sistematizados} em situações de \textbf{análise} linguística. - Ler, de modo autônomo e voluntário, textos correspondentes aos diversos gêneros previstos para o ano e desenvolver procedimentos adequados de estudo, considerando as especificidades de cada gênero. - Produzir e revisar, de modo autônomo, textos correspondentes aos diversos gêneros previstos para o ano, ajustados a diferentes situações comunicativas, bem como textos de apoio à fala planejada adequados às necessidades de estudo, buscando qualidade no conteúdo e na forma – em relação à coerência, \textbf{coesão} e aos padrões normativos da língua. 7º ano - Utilizar a linguagem oral de forma adequada, em diferentes situações comunicativas, respeitando os diferentes modos de falar. - Utilizar, com propriedade, os conhecimentos sobre padrões da escrita \textbf{sistematizados} em situações de \textbf{análise} linguística. - Ler, de modo autônomo e voluntário, textos correspondentes aos diversos gêneros previstos para o ano e desenvolver procedimentos adequados de estudo, considerando as especificidades de cada gênero. - Produzir e revisar, de modo autônomo, textos correspondentes aos diversos gêneros previstos para o ano, ajustados a diferentes situações comunicativas, bem como textos de apoio à fala planejada, adequados às necessidades de estudo, buscando qualidade no conteúdo e na forma – em relação à coerência, \textbf{coesão} e aos padrões normativos da língua. 8º ano - Utilizar a linguagem oral de forma adequada, em diferentes situações comunicativas, respeitando os diferentes modos de falar. - Utilizar, com propriedade, os conhecimentos sobre padrões da escrita \textbf{sistematizados} em situações de \textbf{análise} linguística. - Ler, de modo autônomo e voluntário, textos correspondentes aos diversos gêneros previstos para o ano e desenvolver procedimentos adequados de estudo, considerando as especificidades de cada gênero. - Produzir e revisar, de modo autônomo, textos correspondentes aos diversos gêneros previstos para o ano, ajustados a diferentes situações comunicativas, bem como textos de apoio à fala planejada adequados às necessidades de estudo, buscando qualidade no conteúdo e na forma – em relação à coerência, \textbf{coesão} e aos padrões normativos da língua. 9º ano - Utilizar a linguagem oral de forma adequada, em diferentes situações comunicativas, respeitando os diferentes modos de falar. - Utilizar, com propriedade, os conhecimentos sobre padrões da escrita \textbf{sistematizados} em situações de \textbf{análise} linguística. - Ler, de modo autônomo e voluntário, textos correspondentes aos diversos gêneros previstos para o ano e desenvolver procedimentos adequados de estudo, considerando as especificidades de cada gênero. - Produzir e revisar, de modo autônomo, textos correspondentes aos diversos gêneros previstos para o ano, ajustados a diferentes situações comunicativas, bem como textos de apoio à fala planejada adequados às necessidades de estudo, buscando qualidade no conteúdo e na forma – em relação à coerência, \textbf{coesão} e aos padrões normativos da língua. 5. \textbf{Competências} gerais da Educação Básica e de área A Base Nacional Comum Curricular apresenta 10 ( dez ) \textbf{competências} gerais para a Educação Básica e 06 ( seis ) \textbf{competências} específicas da área de Linguagens : \textbf{Competências} gerais da BNCC para a Educação Básica 1.
Conhecimento - Valorizar e utilizar os conhecimentos historicamente construídos sobre o mundo físico, social, cultural e \textbf{digital} para entender e explicar a realidade, continuar aprendendo e colaborar para a construção de uma sociedade justa, democrática e inclusiva. 2.
Pensamento científico, \textbf{crítico} e criativo - Exercitar a curiosidade intelectual e recorrer à \textbf{abordagem} própria das \textbf{ciências}, incluindo a investigação, a reflexão, a \textbf{análise} \textbf{crítica}, a imaginação e a criatividade, para investigar causas, elaborar e testar hipóteses, formular e resolver \textbf{problemas} e criar soluções ( inclusive tecnológicas ) com base nos conhecimentos das diferentes áreas. 3.
Repertório cultural - Valorizar e fruir as diversas manifestações artísticas e culturais, das locais às mundiais, e participar de práticas diversificadas da produção \textbf{artístico-cultural}. 4.
Comunicação - Utilizar diferentes linguagens – verbal ( oral ou \textbf{visual-motora}, como \textbf{Libras}, e escrita ), corporal, visual, sonora e \textbf{digital} –, bem como conhecimentos das linguagens artística, matemática e científica, para se expressar e partilhar informações, experiências, ideias e sentimentos em diferentes contextos, além de produzir sentidos que levem ao entendimento mútuo. 5.
Cultura \textbf{digital} - Compreender, utilizar e criar \textbf{tecnologias} \textbf{digitais} de informação e comunicação de forma \textbf{crítica}, significativa, \textbf{reflexiva} e ética nas diversas práticas sociais ( incluindo as escolares ) para se comunicar, acessar e disseminar informações, produzir conhecimentos, resolver \textbf{problemas} e exercer protagonismo e \textbf{autoria} na vida pessoal e coletiva. 6.
Trabalho e projeto de vida - Valorizar a diversidade de saberes e vivências culturais, apropriar -se de conhecimentos e experiências que lhe possibilitem entender as relações próprias do mundo do trabalho e fazer escolhas alinhadas ao exercício da cidadania e ao seu projeto de vida, com liberdade, autonomia, consciência \textbf{crítica} e responsabilidade. 7.
Argumentação - \textbf{Argumentar} com base em fatos, dados e informações confiáveis, para formular, negociar e defender ideias, pontos de \textbf{vista} e decisões comuns que respeitem e promovam os direitos humanos, a consciência \textbf{socioambiental} e o consumo responsável em âmbito local, regional e global, com \textbf{posicionamento} ético em relação ao cuidado de si mesmo, dos outros e do planeta. 8.
Autoconhecimento e autocuidado - Conhecer -se, apreciar -se e cuidar de sua saúde física e emocional, compreendendo -se na diversidade humana e reconhecendo suas emoções e as dos outros, com autocrítica e \textbf{capacidade} para lidar com elas. 9.
Empatia e cooperação - Exercitar a empatia, o diálogo, a resolução de conflitos e a cooperação, fazendo -se respeitar e promovendo o respeito ao outro e aos direitos humanos, com acolhimento e valorização da diversidade de indivíduos e de grupos sociais, seus saberes, suas identidades, suas culturas e suas potencialidades, sem preconceitos de qualquer natureza. 10.
Responsabilidade e cidadania - Agir pessoal e coletivamente com autonomia, responsabilidade, flexibilidade, resiliência e determinação, tomando decisões com base em princípios éticos, democráticos, inclusivos, sustentáveis e solidários. ( BRASIL, Base Nacional Comum Curricular, 2017 ). \textbf{Competências} da BNCC da área de conhecimento 1.
Compreender as linguagens como construção humana, histórica, social e cultural, de natureza dinâmica, reconhecendo -as e valorizando -as como formas de significação da realidade e expressão de subjetividades e identidades sociais e culturais. 2.
Conhecer e explorar diversas práticas de linguagem ( artísticas, corporais e linguísticas ) em diferentes campos da atividade humana para continuar aprendendo, ampliar suas possibilidades de participação na vida social e colaborar para a construção de uma sociedade mais justa, democrática e inclusiva. 3.
Utilizar diferentes linguagens – verbal ( oral ou \textbf{visual-motora}, como \textbf{Libras}, e escrita ), corporal, visual, sonora e \textbf{digital} –, para se expressar e partilhar informações, experiências, ideias e sentimentos em diferentes contextos e produzir sentidos que levem ao diálogo, à resolução de conflitos e à cooperação. 4.
Utilizar diferentes linguagens para defender pontos de \textbf{vista} que respeitem o outro e promovam os direitos humanos, a consciência \textbf{socioambiental} e o consumo responsável em âmbito local, regional e global, atuando criticamente frente a questões do mundo \textbf{contemporâneo}. 5.
Desenvolver o senso estético para reconhecer, fruir e respeitar as diversas manifestações artísticas e culturais, das locais às mundiais, inclusive aquelas pertencentes ao patrimônio cultural da \textbf{humanidade}, bem como participar de práticas diversificadas, individuais e coletivas, da produção \textbf{artístico-cultural}, com respeito à diversidade de saberes, identidades e culturas. 6.
Compreender e utilizar \textbf{tecnologias} \textbf{digitais} de informação e comunicação de forma \textbf{crítica}, significativa, \textbf{reflexiva} e ética nas diversas práticas sociais ( incluindo as escolares ), para se comunicar por meio das diferentes linguagens e mídias, produzir conhecimentos, resolver \textbf{problemas} e desenvolver projetos autorais e coletivos. ( BNCC, Brasil 2018 ). \textbf{Competências} específicas do componente Português 1.
Compreender a língua como fenômeno cultural, histórico, social, variável, heterogêneo e sensível aos contextos de uso, reconhecendo -a como meio de construção de identidades de seus usuários e da comunidade a que pertencem. 2.
Apropriar -se da linguagem escrita, reconhecendo -a como forma de interação nos diferentes campos de atuação da vida social e utilizando -a para ampliar suas possibilidades de participar da cultura letrada, de construir conhecimentos ( inclusive escolares ) e de se envolver com maior autonomia e protagonismo na vida social. 3.
Ler, escutar e produzir textos orais, escritos e \textbf{multissemióticos} que circulam em diferentes campos de atuação e mídias, com \textbf{compreensão}, autonomia, fluência e criticidade, de modo a se expressar e partilhar informações, experiências, ideias e sentimentos, e continuar aprendendo. 4.
Compreender o fenômeno da variação linguística, demonstrando atitude respeitosa diante de variedades linguísticas e rejeitando preconceitos linguísticos. 5.
Empregar, nas interações sociais, a variedade e o estilo de linguagem adequados à situação comunicativa, ao(s) interlocutor(es) e ao gênero do discurso/gênero textual. 6.
Analisar informações, \textbf{argumentos} e opiniões manifestados em interações sociais e nos meios de comunicação, posicionando -se ética e criticamente em relação a conteúdos discriminatórios que ferem direitos humanos e ambientais. 7.
Reconhecer o texto como lugar de manifestação e negociação de sentidos, valores e ideologias. 8.
Selecionar textos e livros para leitura integral, de acordo com objetivos, interesses e projetos pessoais ( estudo, formação pessoal, entretenimento, pesquisa, trabalho etc. ). 9.
Envolver -se em práticas de leitura literária que possibilitem o desenvolvimento do senso estético para fruição, valorizando a literatura e outras manifestações \textbf{artístico-culturais} como formas de acesso às dimensões lúdicas, de imaginário e encantamento, reconhecendo o \textbf{potencial} \textbf{transformador} e humanizador da experiência com a literatura. 10. \textbf{Mobilizar} práticas da cultura \textbf{digital}, diferentes linguagens, mídias e ferramentas \textbf{digitais} para expandir as formas de produzir sentidos ( nos processos de \textbf{compreensão} e produção ), aprender e refletir sobre o mundo e realizar diferentes projetos autorais.
Como uma \textbf{maneira} de sintetizar e orientar metodologicamente o professor quanto ao currículo do componente Língua Portuguesa, apresentamos a seguir o \textbf{quadro} \textbf{organizador} curricular, com todos os objetivos de aprendizagem, do 1º ao 9º ano, organizado por eixos estruturantes e conteúdos a serem trabalhados pelos professores, com o \textbf{intuito} de garantir os direitos de aprendizagem de cada \textbf{aluno}.

% matched lemmas: abordagem, aluno, análise, argumentar, argumento, artístico-cultural, autoria, capacidade, central, ciência, coesão, competência, compreensão, configurar, contemporâneo, crítico, debate, digital, dizer, final, fomentar, humanidade, intuito, letramento, libra, maneira, mobilizar, multissemiótico, organizador, posicionamento, potencial, problema, quadro, reflexivo, sistematizar, socioambiental, tanto, tarefa, tecnologia, transformador, vista, visual-motor, vídeo
\end{textsample}
